\section{Основные модули системы}

\newcommand{\algorithmheading}[1]{\par{\leftskip=\parindent\noindent\refstepcounter{subsubsection}\normalfont\fontsize{14pt}{16pt}\selectfont\thesubsubsection\quad #1\par}}

\subsection{Модуль аутентификации и управления доступом}

\headingtextsep

Модуль аутентификации и управления доступом обеспечивает вход пользователей в систему, идентификацию участников рабочих пространств и проверку прав на выполнение операций. Для пользовательского входа применяется внешний провайдер идентификации, а внутренняя модель доступа строится независимо от конкретного провайдера.

В рамках данного модуля поддерживаются регистрация и сопровождение сущности участника рабочего пространства, назначение ролей, приглашение новых пользователей, отключение доступа и передача прав владения рабочим пространством. На уровне ролей должны различаться как минимум владелец, редактор и наблюдатель.

Модуль также отвечает за разделение прав между рабочими разделами проекта, настройками рабочего пространства и системными функциями. Для веб-доступа используются защищённые пользовательские сессии, а изменяющие операции должны проходить с учётом проверок авторизации и защиты от несанкционированных действий.

\textheadingsep
\subsection{Модуль рабочих пространств и проектов}

\headingtextsep

Модуль рабочих пространств и проектов задаёт верхнеуровневую организацию данных системы. Он обеспечивает создание нескольких рабочих пространств, ведение проектов внутри них и логическое разделение данных между разными командами или инициативами.

Рабочее пространство выступает общей административной и организационной границей, внутри которой существуют участники, проекты, агенты и интеграции. Проект является основной рабочей единицей, в пределах которой ведутся задачи, документация, заметки, поиск и история изменений.

Данный модуль должен поддерживать навигацию между рабочими пространствами и проектами, хранение базовых атрибутов проектов, их жизненный цикл и разграничение доступа к проектным сущностям. Все основные объекты системы должны быть связаны с конкретным рабочим пространством и проектом, если они относятся к проектной деятельности.

\textheadingsep
\subsection{Модуль управления задачами}

\headingtextsep

Модуль управления задачами предназначен для постановки, декомпозиции, исполнения и контроля работ. В его составе используются сущности задачи, группы задач и зависимости между задачами, что позволяет описывать как отдельные рабочие элементы, так и крупные логические блоки работ.

Для задач должны поддерживаться создание, редактирование, архивирование, изменение статуса, приоритета, сроков и исполнителя. Исполнителем может быть как участник рабочего пространства, так и программный агент. Для крупных задач и направлений должна поддерживаться группировка в инициативы, эпики и иные структурные объединения.

Модуль должен обеспечивать хранение и обработку зависимостей между задачами. На их основе вычисляются блокировки и определяется порядок выполнения работ. Представление задач в виде списка, доски, иерархии и других форм должно строиться на единой модели данных, а не на независимых реализациях.

\textheadingsep
\subsection{Модуль документации, заметок и файлов}

\headingtextsep

Модуль документации, заметок и файлов отвечает за накопление и сопровождение проектных знаний. В его состав входят страницы документации, быстрые заметки и файловые ресурсы, которые связаны с задачами, группами задач, проектами и результатами работы агентов.

Документация должна храниться в виде страниц с содержимым в формате Markdown. Для неё предусматриваются древовидная структура, внутренние ссылки, поддержка изображений и других вложений, а также возможность связывания со связанными рабочими объектами системы.

Заметки используются для фиксации промежуточного контекста, рабочих наблюдений, решений и итогов выполнения. Файловые ресурсы хранятся отдельно от самих документов: в системе сохраняются их метаданные и связи, а бинарное содержимое размещается во внешнем файловом backend-слое через абстракцию хранения.

\textheadingsep
\subsection{Модуль поиска, интеграций и аудита}

\headingtextsep

Модуль поиска, интеграций и аудита объединяет сервисные функции, обеспечивающие связанность данных и прослеживаемость изменений. Его задача состоит в том, чтобы пользователь или агент могли быстро находить нужную информацию, видеть связанные внешние артефакты и восстанавливать историю значимых действий.

Подсистема поиска должна поддерживать полнотекстовый и семантический поиск по задачам, документам, заметкам, именам и метаданным файлов, а также по результатам агентной активности. Поисковая выдача должна формироваться на основе объединения нескольких механизмов индексации и позволять фильтрацию по типу сущности, проекту и другим признакам.

Подсистема интеграций обеспечивает подключение внешних сервисов, прежде всего GitHub. Она должна поддерживать связывание внутренних сущностей с issue, pull request, commit и branch, а также хранение параметров подключения и состояния внешних интеграций.

Подсистема аудита предназначена для журналирования значимых событий. В журнал должны попадать бизнес-операции, административные действия, события входа, изменения прав, выпуск и отзыв агентных ключей, а также действия, инициированные внешними интеграциями.

\textheadingsep
\subsection{Модуль агентного взаимодействия}

\headingtextsep

Модуль агентного взаимодействия обеспечивает работу программных агентов с системой в контролируемом режиме. В нём описываются сущности агента, его учётных данных, сессий работы и связей между агентной сессией и задачами, с которыми агент взаимодействовал.

Доступ агентов организуется не напрямую к базе данных, а через отдельную MCP-утилиту, работающую поверх основного API. Это позволяет ограничивать набор разрешённых операций, применять единые правила авторизации и фиксировать следы действий агента в журнале.

Для агентных учётных данных должны поддерживаться выпуск, отзыв, ротация, ограничение области действия и проверка статуса. Каждый ключ должен храниться в защищённом виде и предоставлять только явно разрешённые операции, например чтение задач и документов, создание заметок или ограниченное изменение статуса задачи.

Работа агента должна фиксироваться через сессии, связанные с задачами и иными сущностями проекта. Это позволяет прослеживать, какие задачи агент использовал как основной контекст, какие объекты изменял и какие результаты сформировал в ходе выполнения.

\textheadingsep
\subsection{Блок-схемы алгоритмов}

\headingtextsep

Рассмотрим некоторые основные алгоритмы работы системы, связанные с типовыми операциями пользователя и обработкой проектных данных. Они отражают базовые действия, которые можно далее представить в виде компактных блок-схем.

\algorithmheading{Алгоритм создания задачи}

Алгоритм создания задачи начинается с ввода пользователем основных атрибутов: наименования, описания, срока выполнения, приоритета и исполнителя. После получения данных система проверяет заполнение обязательных полей и корректность формата введённых значений. Если проверка проходит успешно, создаётся новая запись задачи, после чего она сохраняется в базе проектных данных. На завершающем шаге система фиксирует событие создания в журнале аудита и возвращает пользователю карточку новой задачи. Блок-схема алгоритма приведена на рисунке А1.

\algorithmheading{Алгоритм добавления заметки}

Алгоритм добавления заметки используется для фиксации промежуточного контекста, решения или результата работы. Пользователь либо агент передаёт текст заметки и указывает объект, с которым она должна быть связана. Далее система проверяет, что содержимое заметки не является пустым и что у автора есть право на создание записи в пределах выбранного проекта. После успешной проверки заметка сохраняется, связывается с нужной сущностью и становится доступной для дальнейшего просмотра в интерфейсе. Блок-схема алгоритма приведена на рисунке А2.

\algorithmheading{Алгоритм выполнения поискового запроса}

Алгоритм выполнения поискового запроса начинается с ввода пользователем поисковой строки и, при необходимости, параметров фильтрации. Затем система определяет область поиска и обращается к индексированным данным по задачам, документам, заметкам и файловым ресурсам. Найденные результаты объединяются, упорядочиваются по релевантности и передаются в интерфейс в виде списка подходящих объектов. Если совпадения отсутствуют, система формирует пустой результат и уведомляет пользователя об отсутствии найденных данных. Блок-схема алгоритма приведена на рисунке А3.

\algorithmheading{Алгоритм загрузки файлового ресурса}

Алгоритм загрузки файла начинается с выбора пользователем нужного ресурса и передачи его в систему вместе с основными метаданными. После этого система проверяет допустимость размера и типа файла, а затем сохраняет бинарное содержимое во внешнем хранилище. После успешной загрузки в базе данных фиксируются метаданные файла, ключ хранения и связи с проектом либо другими сущностями. На заключительном шаге система возвращает сведения о загруженном ресурсе, чтобы он мог быть использован в документации, заметках или других разделах системы. Блок-схема алгоритма приведена на рисунке А4.